\newpage
\section{误差分析}

\subsection{定价模型误差}
\begin{enumerate}
    \item[(1)] \textbf{价格弹性估计误差}：价格弹性系数$\eta，标准误差约0.15。实际弹性因用户异质性而异——商务用户弹性小（$\eta\approx-0.3$），私家车主弹性大（$\eta\approx-0.8$）。
    \item[(2)] \textbf{竞争效应简化}：模型假设相邻充电站竞争效应对称，实际中用户品牌偏好和APP推荐可能导致非对称竞争。
\end{enumerate}
灵敏度分析表明，价格弹性参数$\pm 20\%，最优定价变化不超过5\%，模型具有较好的鲁棒性。

\subsection{预测模型误差}
使用时段聚类的方法存在\textbf{类内变异}——同一日型内的充电行为仍存在差异。聚类准确率89.4\%，意味着约10\%的样本被错误分类，导致对应时段预测偏差。建议引入更多特征（如天气、节假日、促销活动）细化日型分类，降低聚类误差。